\vspace{-0.2in}
\section{Related Work}
\vspace{-0.05in}

\textbf{LLM Safety:} Large language models are typically trained on broad, uncurated datasets, exposing them to harmful content and behaviors. To mitigate these risks, LLMs are often aligned to human preferences via reinforcement learning~\citep{christiano2017deep,bai2022training} or instruction tuning~\citep{ouyang2022training}, and are sometimes paired with content moderation modules~\citep{inan2023llama,zeng2024shieldgemma,han2024wildguard}. Despite these efforts, recent studies have shown that both alignment and moderation can be circumvented by adversarial prompts or jailbreak techniques~\citep{carlini2023aligned,chao2024jailbreakbench,wei2023jailbroken,zou2023universal,liu2023autodan,liu2024calibration}. Most prior work has focused on LLM safety in isolation or in conversational settings. In contrast, our work evaluates LLM safety in the context of both single- and multi-agent systems, where agents may autonomously invoke tools and interact with complex environments.
    
    \textbf{Attacks on LLM-based Agents:}
    Misalignment or jailbreaking of LLMs manifests in the form of toxic content or spread of misinformation in conversational applications. The state-of-the-art LLMs are also capable of using tools and writing code. Unfortunately, safety aligned LLMs can be easily jailbroken in agentic settings leading to scenarios such as generating and executing malicious code~\citep{guo2024redcode}, harmful browser interactions~\cite{kumar2024refusal,tur2025safearena} and multi-step agent misuse~\citep{andriushchenko2024agentharm}. In addition to user prompts through which jailbreaking attacks can be launched, agents are also susceptible to attacks through malicious tool outputs~\citep{debenedetti2024agentdojo,zhang2024agent,zhan2024injecagent,ruan2023identifying} and memory or knowledge-base poisoning~\citep{zhang2024agent,chen2024agentpoison} even when the user prompts are benign. Many frontier LLMs are capable of handling multimodal inputs and are prone to misuse through malicious prompts~\citep{tur2025safearena} and image-based adversarial attacks~\citep{aichberger2025attacking,wudissecting}.
    
    Multi-agent systems introduce additional risks, such as the propagation of malicious prompts between agents~\citep{lee2024prompt}, attacks that exploit agent specialization and collaboration~\citep{tian2023evil,amayuelas2024multiagent}, and vulnerabilities to rogue or compromised agents~\citep{barbi2025preventing}. However, most existing studies focus on specific domains or agent types and use custom, non-comparable evaluation protocols.
    
    \textbf{Agentic Defenses:}
    Safety aligned LLMs and content moderation can be applied for defending agents. However, due to the dynamic nature of agents, another class of defense, based on safety agents is emerging. Given a safety specification, GuardAgent~\citep{xiang2024guardagent} synthesizes a plan and executable code to guard an agent against violations of the specification. AGRail~\citep{luo2025agrail} synthesizes adaptive safety checks based on task-specific requirements, whereas ShieldAgent~\citep{chen2025shieldagent} generates shields that employ probabilistic logical reasoning to monitor action trajectories generates by agents. CaMeL~\citep{debenedetti2025defeating} extracts control and data flow from prompts and uses a custom Python interpreter to enforce fine-grained security policies so that untrusted data cannot impact agent's control flow.
    
    For MAS, AutoDefense~\citep{zeng2024autodefense} filters LLM responses to prevent jailbreak attacks.
    \citet{huang2024resilience} propose a mechanism to improve resilience of multi-agent systems against faulty or malicious agents by allowing agents to challenge messages received from other agents and an extra agent that can inspect and correct messages. To prevent spread of malicious instructions through multi-hop message passing, \citet{peigne2025multi} propose safety instructions and seeding agent memory with examples of safe handling of malicious inputs.
    